% !TEX program = xelatex
% Copyright 2026 Jiayi
% GitHub: @langonginc
% 全国大学生数学建模竞赛社区论文模板（2026 年修订稿）
% 电子版论文第一页必须是摘要页；不要加入承诺书、编号专用页或身份信息。
\documentclass[UTF8,a4paper,zihao=-4]{ctexart}
\usepackage{cumcm2026}

\title{C题示例：数据驱动的预测、评价与优化}
\mcmkeywords{数据预处理；机器学习；预测评价；优化决策}

\begin{document}

% 摘要页即电子版第一页，页码从本页的 1 开始；摘要原则上不超过一页。
\pagenumbering{arabic}
\begin{mcmabstract}
% 用 3--5 句话依次概括研究对象、主要方法、关键结果与结论；正式提交前删除本注释。
本文针对【研究对象与核心任务】，依次完成数据预处理、特征分析、模型构建与结果检验。针对问题一，采用【方法】得到【关键结果】；针对问题二和问题三，分别建立【模型】并给出【结论】。敏感性分析表明【稳定性结论】。本文方法可为【应用场景】提供决策依据。
\end{mcmabstract}

\clearpage
% 正文从第 2 页开始；不要生成目录，正文不得超过 30 页。
\section{问题重述}

\subsection{问题背景}
% 用自己的语言概括题目背景、数据来源与现实意义，避免照抄题面。
在此概括研究背景、业务场景及建模目标，并说明问题属于预测、评价、分类或优化中的哪一类。

\subsection{问题提出}
% 将题目任务改写为彼此衔接、可检验的建模问题。
\begin{problemlist}
  \item 在此说明问题一的输入、输出与评价目标；
  \item 在此说明问题二需要建立的模型及约束；
  \item 在此说明问题三的预测、优化或推广任务。
\end{problemlist}

\section{问题分析}

\subsection{问题一分析}
问题一的核心是建立可靠的数据基础。先检查取值范围、重复记录和缺失比例，再根据指标含义选择中位数填补、截尾或保留缺失标记等方法，最后分别处理效益型和成本型指标。

\subsection{问题二分析}
问题二需要把多维指标映射为可解释的评价结果。可结合题意、专家判断或客观赋权确定权重，再通过线性加权、TOPSIS 等方法计算综合得分，并把资源配置写成带预算和容量约束的优化问题。

\subsection{问题三分析}
问题三关注结论是否依赖某一组参数。可对权重、预算或观测误差进行扰动，比较排序、目标函数与关键决策变量的变化，并据此划定模型的稳定适用区间。

\section{模型假设}

% 假设应有依据、可解释，并在后文说明其影响；不要罗列与模型无关的常识。
\begin{enumerate}
  \item 假设【数据记录或抽样过程】能够代表研究对象的主要特征；
  \item 假设【研究期内的关键关系】保持相对稳定；
  \item 假设【未观测因素】可通过误差项或敏感性分析进行评估。
\end{enumerate}

\section{符号说明}

% symboltable 为可跨页三线表；可选参数为表题，第三列填写单位或“--”。
\begin{symboltable}[主要符号说明]
  \label{tab:symbols}
  $x_{ij}$ & 第 $i$ 个样本在第 $j$ 个特征上的原始观测值 & -- \\
  $z_{ij}$ & 数据预处理后的特征值 & -- \\
  $\boldsymbol{\theta}$ & 待估计的模型参数向量 & -- \\
  $\hat{y}_i$ & 第 $i$ 个样本的预测值 & 题目给定单位 \\
  $L$ & 模型的损失函数或评价指标 & -- \\
\end{symboltable}

\section{数据预处理与探索性分析}

对效益型指标采用极差标准化：
\begin{equation}
  z_{ij}=\frac{x_{ij}-\min_i x_{ij}}{\max_i x_{ij}-\min_i x_{ij}}.
  \label{eq:benefit}
\end{equation}
成本型指标则将式~\ref{eq:benefit} 的分子改为 $\max_i x_{ij}-x_{ij}$。当某指标极差为零时，将其标准化值统一记为零，避免除零错误，并在结果解释中说明该指标没有区分度。

探索性分析至少应报告样本量、缺失比例、极端值和主要相关关系。若使用网上自主查询的数据，应记录来源、访问日期和清洗步骤，并将实际使用的数据文件纳入支撑材料；赛题提供的原始数据不需要重复提交。

\section{问题一模型建立与求解}

\subsection{问题分析}
% 明确问题一的输入、输出、评价指标和建模难点。
在此给出问题一的技术路线。图~\ref{fig:problem1-flow} 可替换为实际算法流程图。
\figureplaceholder{问题一建模流程图}{fig:problem1-flow}

\subsection{数据处理}
% 说明本问题实际使用的样本、特征、训练集/测试集划分及防止数据泄漏的措施。
在此引用第五节的数据预处理结果，并说明问题一的专用处理步骤。

\subsection{模型建立}
% 给出变量、目标函数、约束和参数含义，而不是只列算法名称。
用占位表达式表示预测或评价模型：
\begin{equation}
  \hat{y}_i=f(\boldsymbol{x}_i;\boldsymbol{\theta}),
  \qquad
  \boldsymbol{\theta}^{*}=\arg\min_{\boldsymbol{\theta}}L(\boldsymbol{\theta}).
  \label{eq:model-placeholder}
\end{equation}

\subsection{模型求解}
% 说明算法步骤、参数设置、软件环境、随机种子和停止条件。
在此描述训练、搜索或优化过程，并把完整可运行程序放入附录与支撑材料。

\subsection{结果分析}
% 报告点估计、误差、置信区间或分类指标，并解释结果的实际含义。
\begin{table}[H]
  \centering
  \caption{问题一结果占位表}
  \label{tab:problem1-result}
  \begin{tabularx}{0.78\textwidth}{cXXX}
    \toprule
    模型 & 训练集指标 & 验证集指标 & 备注 \\
    \midrule
    基准模型 & 【填写】 & 【填写】 & 用于比较 \\
    改进模型 & 【填写】 & 【填写】 & 最终模型 \\
    \bottomrule
  \end{tabularx}
\end{table}

\section{问题二模型建立与求解}

在综合得分基础上，以总效益最大为目标建立资源配置模型：
\begin{equation}
  \max \sum_{i=1}^{n} S_i y_i,
  \qquad
  \text{s.t.}\ \sum_{i=1}^{n}y_i\leq B,\quad 0\leq y_i\leq u_i.
  \label{eq:allocation}
\end{equation}
其中 $u_i$ 为对象 $i$ 的容量上限。若资源只能按整数单位分配，应补充 $y_i\in\mathbb{Z}$；若存在公平性、区域覆盖或风险约束，也应根据题意显式加入模型，而不是在求解后人为调整。

示例程序枚举小规模整数方案，便于模板直接运行。实际问题规模较大时，可替换为线性规划、整数规划或启发式算法，并报告求解器设置与终止条件。

\section{问题三模型建立与求解}

\subsection{问题分析}
% 明确问题三是否涉及未来预测、情景模拟、资源配置或策略优化。
在此说明问题三的决策变量、约束和不确定性来源。

\subsection{模型建立}
% 给出可执行的预测或优化模型，必要时说明与前两问模型的耦合关系。
在此写出目标函数、约束条件及模型假设。

\subsection{模型求解}
% 说明滚动预测、启发式算法、整数规划或其他求解步骤。
在此说明参数设定、求解流程及计算复杂度。

\subsection{结果分析}
% 对比不同情景，报告最优方案、预测区间和关键变量变化。
在此给出问题三的定量结论和可执行建议。

\subsection{模型检验与敏感性分析}

% 本文件承接“八、问题三模型建立与求解”，不单独新增一级标题。
% 可使用留出验证、交叉验证、残差分析、扰动实验或极端情景检验。
在此报告误差指标、稳健性结果和参数变化的临界区间，并说明模型失效条件。

\section{模型评价}

\subsection{模型优点}
% 从准确性、可解释性、稳健性、计算效率和可复现性中选择与本文相关的内容。
在此概括模型相较基准方法的主要优势，并用前文结果支撑判断。

\subsection{模型缺点}
% 说明数据、假设、算法或应用边界上的真实局限，避免空泛表述。
在此说明模型依赖的条件、可能的偏差及尚未解决的问题。

\subsection{模型推广}
% 说明模型可迁移到哪些相似任务，以及迁移时需要调整的数据和参数。
在此给出模型推广场景、必要修改和后续研究方向。


% 2026 官方规定：AI 工具使用声明必须位于参考文献之前。
% 2026 官方规定要求在参考文献之前设置本节，并按真实情况二选一。
% 默认启用“未使用 AI”版本，只是为了让模板开箱可编译；正式提交前必须核实并替换。
\section*{AI工具使用声明}

本参赛队在竞赛过程中未使用任何AI工具。

% 若真实使用了 AI，请删除上面一句，取消下面一句的注释，并如实填写〖〗中的简要用途：
% 本参赛队在竞赛过程中使用了AI工具，主要用于〖简要用途，如语言润色、代码调试等〗，详细使用情况见支撑材料。
%
% 同时必须在支撑材料中提供“AI工具使用详情.pdf”，包括：
% 1. 所用 AI 工具名称、版本或型号；
% 2. 具体使用目的和环节；
% 3. 主要提示方式与使用过程说明（可附典型交互示例）；
% 4. 对 AI 输出的采纳、人工修改和核验的主要情况（语言润色除外）。

% 参考文献单独存放在此文件。
% 在正文中使用 \cite{条目标签} 引用；每个 \bibitem 的标签应唯一。
\begin{thebibliography}{9}
  \bibitem{topsis}
  Hwang, C. L., and Yoon, K. \textit{Multiple Attribute Decision Making: Methods and Applications}. Springer, 1981.
\end{thebibliography}


\mcmappendix
\section{支撑材料文件列表}
% 仅列出实际提交的支撑材料。赛题提供的原始数据无需重复列作自主数据。
\begin{itemize}
  \item \texttt{code/problem1.py}、\texttt{code/problem2.py}、\texttt{code/problem3.py}：完整可运行源程序；
  \item \texttt{code/problem1.m}：Matlab 程序示例；
  \item \texttt{data/external-data.xlsx}：自主查询并实际使用的数据（如有）；
  \item \texttt{results/intermediate-results.xlsx}：支撑结论的必要中间结果（如有）；
  \item \texttt{AI工具使用详情.pdf}：仅在真实使用 AI 工具时按官方规定提供。
\end{itemize}

\section{完整源程序代码}
% 正式提交时，应展示建模所用到的全部完整、可运行程序。
\sourcefile[language=Python]{问题一完整源程序}{code/problem1.py}
\sourcefile[language=Python]{问题二完整源程序}{code/problem2.py}
\sourcefile[language=Python]{问题三完整源程序}{code/problem3.py}
\sourcefile[language=Matlab]{问题一 Matlab 程序示例}{code/problem1.m}

\end{document}
